\chapter{光学：最短时间原理}

\section{光}

这是关于电磁辐射主题诸章中的第一章。我们看到的光，在同为辐射的宽广谱系中只占一小部分，这个谱系的各个部分是以某一个变量的不同数值来区分的。这个变量可以叫做“波长”。当它在可见光谱中变化时，光明显地从红色变到紫色。如果我们要系统地从长波长到短波长来考察这个谱系，那么就要从通常所谓的无线电波开始。工程技术上所能得到的无线电波，其波长很宽，有的甚至比通常广播上所用的波长还长；通常广播用的波长大约为$500~m$。接着是所谓“短波”，即雷达波、毫米波等等。在一个波长范围与另一个波长范围之间并没有实际的界线，因为大自然没有为我们提供明确的边界。与某一波段名称相联系的波长数值只是近似的，因而我们给不同波段所取的名字当然也是近似的。

在通过相当长一段毫米波之后，我们就来到了所谓的红外区，接着进入到可见光谱区。朝波长更短的方向走去，就进入所谓的紫外区。紫外区终止处，就开始出现X光区。但我们不能精确地确定这个分界处在哪里；它大致在$10^{-8}~m$即$10^{-2}\mu m$处。这些是“软”X射线；随着所谓波长这一值的越来越小，接着是一般X射线和硬X射线；再就是$\gamma$射线等等。

在这个宽广的波长范围内，有三个或更多个大致的范围特别令人感兴趣。其中的一个区域满足这样的条件，即所涉及的波长比用来研究它们的仪器的尺度要小得多；而且，用量子论来说，其光子的能量比仪器的能量灵敏度要小，在这些条件下，我们可以用一种叫做几何光学的方法来作粗略的一级近似。另一区域是，如果波长与仪器的尺度可以相比拟（这对可见光虽然难以实现，但对无线电波却比较容易做到），而光子的能量仍小得可以忽略不计，那么只要研究波的行为，就可以得到一种很有用的近似，仍然不必去考虑量子力学。这一方法是建立在电磁辐射经典理论的基础之上的。这一理论将在较后的一章中讨论。下一个，如果我们考察很短的波长，这里波的特性可以忽略，但光子却具有比仪器的灵敏度大得多的能量，那么事情又变得很简单了。这时就得用简单的光子图像，对此我们只能十分粗糙地加以描述。把整个事物统一于一个模型的完整图像，在相当长一段时间内我们将用不到。

在本章中我们的讨论将限于几何光学范围，在这里我们不去顾及光的波长和粒子特征，而这些都将在适当的时候再作解释。我们甚至不必操心去问光是什么，而只是要找出在与有关的尺度相比很大的范围内光的行为如何。所有这些都必须事先加以说明，以便强调这一事实，即我们打算讲述的内容仅仅是一种十分粗糙的近似；这是我们以后必须重新“忘掉”的几章之一。而且我们将会很快忘掉它，因为我们几乎马上就要讲到一个更精确的方法。

虽然几何光学只是一种近似的方法，但它在技术上非常重要，历史上也令人极感兴趣。为了使读者对物理理论或物理概念的发展获得一些知识，我们将把这个课题比别的一些课题更多地从历史上来加以叙述。

首先，光无疑是每个人所熟悉的，而且从无法回忆的时间起就早已熟悉了。现在的问题是，我们凭什么作用看见光？关于这个问题曾经提出过许多理论，但最后可归结到一个，即认为有某种东西进入眼睛——它是从物体上弹出而进入眼睛的。这种想法我们已听了这么久，以致习以为常，因而几乎不可能相信会有某一个十分聪明的人会提出相反的理论——比方说有什么东西由眼睛里射出而感触到了物体。另外一些观察到的重要事实是，当光从一处射到另一处时，如果途中没有障碍物，它就沿直线行进，并且光线之间似乎不会发生相互干扰。这就是说，虽然光在房间里在各个方向上相互交叉地射来射去，但横穿过我们视线的光并不影响从某一物体向我们射来的光。这曾经是惠更斯一度用来反对微粒说的一个最有力的证据。你想，如果光像许许多多飞逝着的箭，那么其他的箭怎么能如此容易地穿过它们？但这样的哲学论据并不具有很大的说服力。人们总是可以说光就是由可以相互贯穿的飞箭组成的！

\section{反射与折射}

下面的讨论为几何光学提供了足够的基本概念，现在我们要稍微深入到定量一些的内容方向。到目前为止，我们只说明了光在两点之间沿直线行进；现在来研究一下光射向各种物质时的行为。最简单的物体是一面镜子，关于镜子的规律是，当光射到镜子时，就不再继续沿着原来的直线行进，而从镜子跳到一条新的直线中，这条直线将随镜子倾角的改变而改变。对于前人来说，他们所面临的问题就是其中两个角之间有什么关系？这是一种很简单的关系，并且很早就已发现了。射向镜子的光的行迹是这样的：两束光的每一束与镜面形成一个夹角，这两个夹角相等（图\ref{figure:入射角等于反射角}）。由于某种原因，通常是从镜面的法线来量度这些角的。因此，所谓的反射定律就是
\begin{equation}
    \label{Eq:I:26:1}
    \theta_i = \theta_r
\end{equation}

这是一个十分简单的命题，但是当光从一种介质进入另一种介质，例如从空气进入水中时，就碰到了一个较为困难的问题；这里与前面一样，我们也看到光不再沿原直线行进（图\ref{figure:当光从一种介质进入另一种介质时发生折射}）。光在水中的路线偏离了它在空气中的路线；如果我们改变$\theta_i$角使之更接近于垂直入射，那么“偏折”角就不大。但是如果使光束的倾角变得很大，那么所偏离的角度就很大。现在的问题是，一个角与另一个角的关系如何？这个问题使前人迷惑了很长一段时间，而且对此他们从未找到过答案！然而，这正是人们在全部希腊物理学中可以发现的用表列出实验结果的少数场合之一。托勒玫把空气中若干个不同的角同与之一一对应的水中的角列成表。表\ref{Table:26:1}列出了以度为单位的空气中的角和在水中量得的相应的角。（一般认为希腊科学家从来不做任何实验。但是，如果不知道正确规律，则除了用实验外，要得到这张表是不可能的。然而必须指出，这些值并不表示对每个角度所作的独立和谨慎的测量，而只是从少数几个测量中用内插法得到的一些数值而已，因为它们与一条抛物线完全相符合。）


\begin{figure}[htbp]
    \centering
    \begin{minipage}[t]{0.45\textwidth}
        \centering
        \includegraphics[width=5.5cm]{Chapter26/入射角等于反射角}    
        \caption{入射角等于反射角}
        \label{figure:入射角等于反射角}
    \end{minipage}
    \qquad
    \begin{minipage}[t]{0.45\textwidth}
        \centering
        \includegraphics[width=5.5cm]{Chapter26/当光从一种介质进入另一种介质时发生折射}
        \caption{当光从一种介质进入另一种介质时发生折射}
        \label{figure:当光从一种介质进入另一种介质时发生折射}
    \end{minipage}
\end{figure}

\begin{table}[!ht]
    \centering
    \caption{}
    \label{Table:26:1}
    \setlength{\tabcolsep}{10mm}{
    \begin{tabular}{|c|c||c|c|}
        \hline
        空气中的角 & 水中的角 & 空气中的角 & 水中的角 \\ \hline
        10$^\circ$    & 8$^\circ$    & 50$^\circ$    & 35$^\circ$   \\ 
        20$^\circ$    & 15.5$^\circ$ & 60$^\circ$    & 40.5$^\circ$ \\ 
        30$^\circ$    & 22.5$^\circ$ & 70$^\circ$    & 45.5$^\circ$ \\ 
        40$^\circ$    & 29$^\circ$   & 80$^\circ$    & 50$^\circ$   \\ \hline
    \end{tabular}
    }
\end{table}

我们首先观察一个效应，接着进行测量并把它列成表；然后试图找出可以把一事物与另一事物联系起来的规律；这是物理定律在发展过程中的重要步骤之一。上面的数值表是公元140年作的，但是直到1621年才有人终于找到了这个联系两个角的规律！这个由荷兰数学家斯涅耳所发现的规律可叙述如下：如果$\theta_i$表示空气中的角，$\theta_r$表示水中的角，则$\theta_i$的正弦等于某一常数乘上$\theta_r$的正弦
\begin{equation}
    \label{Eq:I:26:2}
    \sin\theta_i=n\sin\theta_r
\end{equation}
对于水，常数$n$约为$1.33$。等式（\ref{Eq:I:26:2}）称为斯涅耳定律；它使我们能预言光从空气进入水中时将如何弯折。表\ref{Table:26:2}列出了根据斯涅耳定律得到的空气中的角和水中的角。注意它与托勒玫那张表非常一致。

\begin{table}[!ht]
    \centering
    \caption{}
    \label{Table:26:2}
    \setlength{\tabcolsep}{10mm}{
    \begin{tabular}{|c|c||c|c|}
        \hline
        空气中的角 & 水中的角 & 空气中的角 & 水中的角 \\ \hline
        10$^\circ$    & 7.5$^\circ$    & 50$^\circ$    & 35$^\circ$   \\ 
        20$^\circ$    & 15$^\circ$     & 60$^\circ$    & 40.5$^\circ$ \\ 
        30$^\circ$    & 22$^\circ$     & 70$^\circ$    & 45$^\circ$   \\ 
        40$^\circ$    & 29$^\circ$     & 80$^\circ$    & 48$^\circ$   \\ \hline
    \end{tabular}
    }
\end{table}

\section{费马最短时间原理}

随着科学的进一步发展，我们所要知道的是比仅仅一个公式更多的东西。首先我们作了观察，接着有了通过测量得到的数值，然后有了概括所有这些数值的定律。但是科学的真正光荣在于我们能够找到一种思想方法，使得定律成为明显的。

使得有关光的行为的那个定律成为明显事实的第一个思想方法是费马于1650年左右发现的，它称为最短时间原理或费马原理。他的想法是这样的：在从一点行进到另一点的所有可能的路径中，光走的是需时最短的路径。

首先我们来证明这对镜子的情况是正确的，这个简单原理既包含了直线传播定律又包括了镜面反射定律。这样，我们的理解就加深了！我们来试着求下列问题的解。在图\ref{figure:最短时间原理的说明}中画了$A$、$B$两点和一平面镜$MM'$。哪一条是在最短时间内从$A$走到$B$的路径呢？
\begin{wrapfigure}{l}{0.35\textwidth}
    \centering
    \includegraphics[width=0.33\textwidth]{Chapter26/最短时间原理的说明}
    \caption{最短时间原理的说明}
    \label{figure:最短时间原理的说明}
\end{wrapfigure}
回答是从$A$笔直走到$B$的那条！但是如果我们附加一条规定，即光必须在最短时间内碰到镜面再返回到$B$，那么回答就不会这么容易了。一个方法是沿着$ADB$路径尽快到达镜子然后再走到$B$。当然，这时要走一条很长的路程$DB$。如果让$D$略向右偏移到$E$，虽然稍稍增加了第一段距离，但却大大减少了第二段距离，这样就使总的路程减少，从而使传播的时间也相应地减少。那么怎样来找到需要时间最短的$C$点呢？回答是：可以很巧妙地用几何技巧来找到它。

在$MM'$的另一边作一个人为的点$B'$，使它在镜面$MM'$以下的距离跟B点在镜面以上的距离相同。然后作$EB'$线。由于$BFM$是直角，且$BF=FB'$，所以$EB$等于$EB'$。于是两段距离之和$AE+EB$（光以不变速度传播时，它与光的传播时间成正比）即为两段长度$AE$与$EB'$之和。因此，问题就变为何时这两段长度之和最短？回答很简单：当线段经过$C$点使从$A$到$B'$成为一直线！换句话说，我们必须找出这个点，通过它能直达人为的点，那么这个点就是所要求的正确的点。现在$ACB'$既为一条直线，则角$BCF$等于角$B'CF$，因而也等于角$ACM$。因此，入射角等于反射角的这种说法与光射向镜面沿着需时最短的路径返回到$B$的说法是等效的。这一命题最初由亚历山大的希罗（Hero）所述：光走到镜面然后射向另一点时所取的是尽可能短的一段距离，因而这并不是一个新的理论。但正是这一点启发了费马，使他联想到光的折射也许是在同一基础上进行的。但在折射中，光显然不取最短距离的路径，所以费马试用了光取最短时间的思想。

在我们对折射进行分析之前，应就镜子再说几句话。如果在$B$点有一光源，它朝镜子发光，那么我们将看到自$B$点射到$A$点的光恰好像没有镜子时在$B'$点有一物体把光射到$A$点一样。眼睛当然只觉察到实际射进它的光，所以如果在$B$点有一物体，并有一面镜子，它使光进入眼睛完全好像物体是在$B'$处而把光射入眼睛一样，那么在视觉-大脑系统不太知道其他情况的假定下，它就会将此解释为有一物体在$B'$处。所以，认为有物体位于镜子后面的这种错觉，仅仅是由于这样的事实，即进入眼睛的光实际上完全好像来自镜子后面的一个物体一样（除了镜面上有灰尘，以及我们知道镜子的存在等等这些可在大脑中被校正的因素外）。

现在我们来证明最短时间原理将导出斯涅耳折射定律。但必须就光在水中的速率作一个假定。我们假定光在水中的速率比在空气中小一个因子$n$。

在图\ref{figure:费马原理对折射的说明}中，我们的问题仍然是在最短时间内从$A$走到$B$。为了说明最好的办法并不是沿着直线走这一事实，设想有一个漂亮的少女从船上掉了下去，她正在水中的$B$点向人们呼救。以$X$标出的线表示河岸。我们是在陆地上的$A$点看到了这件事；我们既能跑步又能游泳，但跑得比游得更快。那么我们怎么办？笔直过去吗（是的，毫无疑问）？但只要稍加思索，就会领悟到在陆地上稍微多跑一些路以减少在水中的路程是有利的，因为我们在水中的速度要慢得多（按此推理，我们将会说正确的做法是要十分仔细地算出应该怎么办）。不论如何，我们来证明问题的最终解答是路径$ACB$，并证明这是所有可能的路径中需时最短的一条。如果它是这样，那就意味着若取任何别的路径，需时就要较长。所以，如果我们以所需的时间对$X$点的位置作图，那么将会得到一条如图\ref{figure:C点相应于最短时间，但其附近各点相应于几乎与之相同的时间}所示的曲线，其中对应于$C$点的时间是所有可能的时间中最短的。这就是说如果我们将$X$点移到邻近$C$的各点，在一级近似下时间基本不变，因为曲线底部的斜率为零。所以，我们寻找这条规律的方法，就是设想把位置作很小的移动而要求时间基本不变（当然有一个二级无穷小的改变；对于从$C$点向两旁任何一个方向的移动来说，都应有一正的增量）。因此我们考察邻近的一点$X$，并算出从$A$到$B$沿这两条路径所需的时间各有多长，再将新路径的时间与旧路径的时间相比较。这很容易做到。当距离$XC$很短时，我们当然要求这个时间之差接近于零。先看岸上的路程。如果作垂线$XE$，就可看出这段路程比$AC$短了一段$EC$。我们说因为没有走这一段额外的距离而得到好处。另一方面，在水中，在作了相应的垂线$CF$后，发现必须多走一段额外的距离$XF$，而使我们有所损失。或者说，在时间上，我们赢得了走过距离$EC$所需的时间，但失掉了走过距离$XF$所需的时间。这两段时间必须相等，因为在一级近似下，总的时间应该不变。假定在水中的速率是空气中速率的$1/n$，则必须有
\begin{equation}
    \label{Eq:I:26:3}
    EC=n \cdot XF
\end{equation}

\begin{figure}[htbp]
    \centering
    \begin{minipage}[t]{0.45\textwidth}
        \centering
        \includegraphics[width=5.5cm]{Chapter26/费马原理对折射的说明}    
        \caption{费马原理对折射的说明}
        \label{figure:费马原理对折射的说明}
    \end{minipage}
    \qquad
    \begin{minipage}[t]{0.45\textwidth}
        \centering
        \includegraphics[width=5.5cm]{Chapter26/C点相应于最短时间，但其附近各点相应于几乎与之相同的时间}
        \caption{C点相应于最短时间，但其附近各点相应于几乎与之相同的时间}
        \label{figure:C点相应于最短时间，但其附近各点相应于几乎与之相同的时间}
    \end{minipage}
\end{figure}

可见，当我们选择了正确的点时，就有$XC\sin\angle EXC=nXC\sin\angle XCF$，消去公共斜边$XC$，并注意到
\begin{equation*}
    \angle EXC = \angle ECN = \theta_i ~~~~\text{及}~~~~ \angle XCF = \angle BCN' = \theta_r
\end{equation*}
就有
\begin{equation}
    \label{Eq:I:26:4}
    \sin\theta_i = n\sin\theta_r
\end{equation}
所以，当速率之比为$n$时，为了在最短时间内从一点到达另一点，光应以这样的角入射，使得角$\theta_i$和$\theta_r$的正弦之比等于光在两种介质中的速率之比。

\section{费马原理的应用}

现在我们来讨论最短时间原理的一些有趣的结果。首先是倒易原理。如果从$A$到$B$我们找到了需时最短的路径，那么沿相反方向走的（假定光沿任一方向行进的速率相等）需时最短的路径将是这同一条路径。因此，如果光可以沿一条路径行进，也就可以倒转过来行进。

另一个有趣的例子是一块具有平行平面的玻璃板，它与光束成一角度。当光从$A$点经过玻璃板走到$B$点时（图\ref{figure:当一束光通过一块透明版时，它被偏移了}），并不沿一直线通过，而代之以在玻璃板中使倾角较小的路径，以减少它在玻璃板中所花的时间，虽然这样使它在空气中所花的时间略有增加。光束只是本身平移了一下，因为入射角与出射角是相等的。

\begin{figure}[htbp]
    \centering
    \begin{minipage}[t]{0.45\textwidth}
        \centering
        \includegraphics[width=5.5cm]{Chapter26/当一束光通过一块透明版时，它被偏移了}    
        \caption{当一束光通过一块透明版时，它被偏移了}
        \label{figure:当一束光通过一块透明版时，它被偏移了}
    \end{minipage}
    \qquad
    \begin{minipage}[t]{0.45\textwidth}
        \centering
        \includegraphics[width=5.5cm]{Chapter26/地平线附件，表观太阳比真实太阳高出约}
        \caption{地平线附件，表观太阳比真实太阳高出约$1/2$度}
        \label{figure:地平线附件，表观太阳比真实太阳高出约}
    \end{minipage}
\end{figure}

第三个有趣的现象是这样一个事实：当我们看见落日时，它其实已在地平线以下了！它看起来似乎并不在地平线以下，但实际上确是如此（图\ref{figure:地平线附件，表观太阳比真实太阳高出约}）。地球的大气高处稀薄而底部稠密。光在真空中传播得比空气中快，因而，如果太阳光不沿地平线行进，而从地平线以上以较陡的倾斜度通过稠密区，以尽量减少光在其中行进得慢的这一区域中的路程的话，它就能较快地到达$S$点。当太阳看来要落到地平线以下时，其实它已落到地平线以下好些时候了。这个现象的另一个例子是当人们在炎热的道路上驾车时，常常会见到的海市蜃楼现象。人们在路上见到“水”，但当他到达那里时，却干燥得像沙漠一样！这一现象可说明如下。其实我们真正见到的乃是从路上“反射”上来的天上的光，如图\ref{figure:海市蜃楼现象}所示。来自天空投射到道路上的光，能朝上到达眼睛。为什么呢？紧靠地面之上的空气很热，但越往高处则越冷。热空气比冷空气膨胀得多一些，因而更稀薄一些，这就使光的速率减小得少一些。也就是说，光在热的区域跑得比冷的区域快。这样一来，为了节省时间，光就选定不沿直线直接过来，而取了费时最短的路径，因而暂时走进它跑得快一些的区域。所以，光能沿着曲线行进。作为最短时间原理的另一个重要例子，假定我们设想一种装置，使所有发自一点$P$的光重新汇集到另一点$P'$（如图\ref{figure:光学“暗箱”}）。这当然意味着光能够沿一直线从$P$到达$P'$点。这完全对。但是，我们如何设法使光不仅能沿直线行进到达$P'$点，而且也能使从$P$点朝$Q$点发出的光终止于$P'$点呢？我们要把所有的光引回到所谓的“焦点”。怎么引呢？如果光老是走需时最短的路径，那么它当然不会沿所有其他的路径走。唯一使光能很好地做到也能沿着邻近的一些路径行进的方法，是使这些路径所需的时间恰好相等！否则，光就会选择需时最短的一条路径走。所以造成一个聚焦系统的问题，仅仅在于设计一个器件，使得光沿着所有不同路径走时所花的时间相等。

\begin{figure}[htbp]
    \centering
    \begin{minipage}[t]{0.45\textwidth}
        \centering
        \includegraphics[width=5.5cm]{Chapter26/海市蜃楼现象}    
        \caption{海市蜃楼现象}
        \label{figure:海市蜃楼现象}
    \end{minipage}
    \qquad
    \begin{minipage}[t]{0.45\textwidth}
        \centering
        \includegraphics[width=5.5cm]{Chapter26/光学“暗箱”}
        \caption{光学“暗箱”}
        \label{figure:光学“暗箱”}
    \end{minipage}
\end{figure}

这很容易做到。假定我们有一片玻璃，光在其中走得比空气中慢（图\ref{figure:一种光学聚焦系统}）。现在考察一条在空气中沿着$PQP'$路径走的光线。这是一条比直接从$P$点到$P'$点长的路径，无疑要花较长的时间。但是如果我们插入一片恰当厚度的玻璃（以后我们将算出它有多厚），它或许恰好能补偿因光倾斜着走而多花的时间！在这种情况下，我们就能使笔直通过的光所花的时间与沿着$PQP'$路径所花的时间相同。同样地，如果我们取一条稍为倾斜的光线$PRR'P'$（它不如$PQP'$那么长），对它的补偿虽然不必像笔直的一条那么多，但总得补偿一些。我们竖着放置一片形如图\ref{figure:一种光学聚焦系统}那样的玻璃。利用这种形状，所有从$P$发射的光就会到达$P'$。当然，这对我们来说很熟悉，我们称这种器件为会聚透镜。在下一章中，我们将实际计算为了达到完善的聚焦点，透镜应具有什么形状。

\begin{figure}[htbp]
    \centering
    \begin{minipage}[t]{0.45\textwidth}
        \centering
        \includegraphics[width=5.5cm]{Chapter26/一种光学聚焦系统}    
        \caption{一种光学聚焦系统}
        \label{figure:一种光学聚焦系统}
    \end{minipage}
    \qquad
    \begin{minipage}[t]{0.45\textwidth}
        \centering
        \includegraphics[width=5.5cm]{Chapter26/椭球面镜}
        \caption{椭球面镜}
        \label{figure:椭球面镜}
    \end{minipage}
\end{figure}

再举一个例子：假如我们想要设置几面镜子，使从$P$发出的光总是到达$P'$（图\ref{figure:椭球面镜}）。在任一路径上，光都射到某一面镜子然后返回到$P'$，而所有这些路径所花的时间必须相等。这里光一直在空气中行进，所以时间与距离成正比。于是所有时间都相同的说法跟所有距离都相同的说法完全一样。因而两段距离$r_1$与$r_2$之和必须是常数。椭圆就是具有这样一种性质的曲线，即从两个点到椭圆上任何一点的距离之和为一常数；这样我们就可保证从一个焦点发出的光都到达另一个焦点。

同样的原理适用于聚集来自一颗星的光。巨大的帕洛玛$200~in$望远镜就是用下述原
\begin{wrapfigure}{r}{0.45\textwidth}
    \centering
    \includegraphics[width=0.44\textwidth]{Chapter26/抛物面镜}
    \caption{抛物面镜}
    \label{figure:抛物面镜}
\end{wrapfigure}
理制成的。设想有一颗数$10^9mi$远的星；我们希望能使所有射来的光线都到达焦点。当然我们不能把一直到星的整段光路画出，但仍要核对一下各条光线所需的时间是否相等。我们当然知道，当各条光线到达与光线垂直的某一平面$KK'$时，它们所花的时间是相等的（图\ref{figure:抛物面镜}）。于是所有光线必须在相等的时间内从这里射到镜面而再继续行进到$P'$。那就是说，我们必须找到一条曲线，它具有这种性质，不论X取在哪里，都能使距离之和$XX'+X'P'$为一常数。寻找它的一个简便方法就是将$XX'$线延长到$LL'$平面。现在，如果我们把一条曲线安排得使$A'A''=A'P'$，$B'B''=B'P'$，$C'C''=C'P'$等等，那么这就是我们所要求的曲线，因为这时$AA'+A'P'=AA'+A'A''$显然是常数。因而我们的曲线就是所有与一点和与一直线等距离的点的轨迹。这样的曲线称为抛物线；镜面就被做成抛物线形状。

上述例子说明了这些光学器件赖以设计的原理。严格的曲线可用这样的原理计算得到，即为了能够完全聚焦，所有光线的传播时间必须精确相等，同时也必须小于邻近任何一条路径所花的时间。

我们将在下一章中进一步讨论这些聚焦用的光学器件；现在我们来讨论这个理论的进一步发展。当一种像最短时间原理那样的新的理论发展起来时，我们首先会倾向于说：“嗯，那很漂亮；令人高兴；但问题是，它最终是否会对理解物理有所帮助？”有人会回答：“有，你看我们现在能理解多少东西呀！”另一个则说：“很好，但我也能理解抛物面镜。我只需要这样一条曲线，使它的每一个切面跟两条光线构成相等的角度。我也能画一个透镜，因为射到它上面的任何一条光线都将按斯涅耳定律给定的角度折射。”显然，最短时间的说法和在反射中两角度相等的说法以及在折射中两角度的正弦成比例的说法都是相同的。所以这是否只是一个哲学问题，或者只是一个审美上的问题？两方面都有充分的论据。

然而，一个有效原理的重要性在于它能预言新的东西。

很容易说明费马原理预言了许多新的东西。首先，假定有三种介质——玻璃、水和空气，我们做一个折射实验以测量一种介质对另一种介质的折射率$n$。把空气（1）对水（2）的折射率叫做$n_{12}$；空气（1）对玻璃（3）的折射率叫做$n_{13}$。如果就水对玻璃的折射率进行测量，那么将会找到另一个我们称之为$n_{23}$的折射率。但并无先验的理由可以说明$n_{12}$、$n_{13}$和$n_{23}$之间应有什么联系。然而，根据最短时间的思想，它们之间确实存在着一个确定的关系。折射率$n_{12}$是两个量，即空气中的速率和水中的速率之比；$n_{13}$是空气中的速率和玻璃中的速率之比；$n_{23}$是水中的速率和玻璃中的速率之比。所以如果消去空气中的速率，就得到
\begin{equation}
    \label{Eq:I:26:5}
    n_{23} = \frac{v_2}{v_3} = \frac{v_1/v_3}{v_1/v_2} = \frac{n_{13}}{n_{12}}
\end{equation}
换句话说，我们可以预言一对新的介质的折射率可以从各个介质对空气或真空的折射率中求得。所以，如果对光在所有介质中的速率进行测量，并由此对每一种介质得到一个单一的数，即它相对于真空的折射率$n_i$（例如$n_1$是真空中速率与空气中速率之比，等等），那么我们的公式就变得非常简单。于是任意两种介质$i$和$j$的折射率为
\begin{equation}
    \label{Eq:I:26:6}
    n_{ij} = \frac{v_i}{v_j} = \frac{n_j}{n_i} 
\end{equation}
只用斯涅耳定律，就不存在作这类预言的基础\footnote{虽然在作一个附加假定的条件下它可以推导出来。这个附加假定是：在一种物质的表面上加一层另一种物质，不会改变光在后一种物质中最终的折射角。}。但这一预言当然成立。关系式（\ref{Eq:I:26:5}）很早就已知道，而且是最短时间原理的一个有力证据。

最短时间原理的另一证据即另一预言是，如果我们测量光在水中的速率，它定将比空气中的慢。这完全是另一种类型的预言。它是一个出色的预言，因为到此为止所有我们测得的量都是角度；而这里我们有了一个与观察十分不同的理论上的预言，费马则是由这些观察引出他的最短时间这个概念的。结果证明，光在水中的速率的确比空气中的速率慢，两者之比恰好等于折射率！

\section{费马原理的更精确表述}

实际上，我们必须把最短时间原理表述得更正确一些。在前面，它并没有得到正确的表述。称它为最短时间原理是不正确的，我们只是为了方便起见，一直沿用着这个不正确的表述，但现在我们必须考虑正确的表述应该是怎样的。假如我们有如图\ref{figure:最短时间原理的说明}所示的一面镜子。是什么促使光一定要跑到镜面上来？最短时间的路径显然是$AB$。所以有些人也许会说：“有时这是一条取时最长的路径。”但这里的时间并不是极大值，因为一条弯曲的路径需要的时间肯定要更长！正确的表述如下：在某一特定路径上行进的光具有这样的一种性质，那就是，如果我们不论用何种方式使光路作微小改变（比如说移动百分之一），例如改变它射到镜面上的路线的位置，或改变曲线的形状，或任何别的方式，都不会有时间的一级变化；而只有时间的二级变化。换句话说，这个原理是：光取的是这样一条路径，在它邻近有许多取时几乎与它完全相同的其他路径。

下面是最短时间原理的另一个困难，它也是那些不喜欢这种理论的人所永远不能忍受的一个困难。利用斯涅耳理论我们能够“理解”光。当光在行进中遇见一个表面时，会因为它与表面有相互作用而发生弯曲。光从一点跑到另一点，再到又一点等等的因果关系的思想是容易理解的。但最短时间原理却是关于自然界行为方式的完全另一种哲学原理。它不是把我们做了某件事后，又发生另一件事等等说成是因果关系，而是这样说：我们建立起一种装置，由光束判断哪一条是需时最短或最终的路径，从而选定由这条路径走。但它做了什么，又是如何找出这条路径的？它是否试探了邻近各条路径，
\begin{wrapfigure}{r}{0.45\textwidth}
    \centering
    \includegraphics[width=0.44\textwidth]{Chapter26/无线电波通过一个狭缝}
    \caption{无线电波通过一个狭缝}
    \label{figure:无线电波通过一个狭缝}
\end{wrapfigure}
并对它们进行了相互核对？回答是：是的，在某种程度上，它的确是这样做的。这个特征在几何光学中当然是无法知道的，但却包含在波长的概念之中；波长大致告诉我们，光必须在离开多远的地方“试探”这条路径以便进行核对。要在大尺度内用光来演示这一事实是困难的，因为光的波长实在太短。但对于无线电波，例如3厘米波，进行核对的距离就比较长了。如果我们有一个无线电波源，一个探测器和一条缝，如图\ref{figure:无线电波通过一个狭缝}所示，波当然从$S$跑到$D$，因为这是一条直线，而且当我们把缝关小时仍然如此——波仍沿此直线走。但如果把探测器移到旁边的$D'$点，波将不会通过宽缝从$S$跑到$D'$，因为它们核对了邻近的几条路径，并且说：“不，我的朋友，那些路径都对应于不同的时间。”但是如果我们把缝关小到使之成一条很窄的缝隙，以阻止辐射波核对各条路径，那么只有一条路径可走，辐射波也就只好随着它走！用一条狭缝比起用宽缝来有更多的辐射能到达$D'$点！

用光也能作同样的实验，但要在大尺度上来演示是困难的。其效应可在下述简单情况下看到。找一小而亮的光源，比如说远处街灯的一个非乳白灯泡，或太阳在弯曲的汽车缓冲器上的反射光。将两只手指放在一只眼睛前面，以便通过指间狭缝进行观察，然后慢慢地并拢手指使光减少到接近于零。此时，你将看到先前是一个小点的光源像变得相当长，甚至拉成一条长线。理由是手指靠得很近，原来认为沿直线过来的光散开成一个角度，因此当它进入眼睛时就从几个方向跑了进来。如果你很仔细的话，还会看到侧向极大，即许多与缝隙边缘平行的条纹，而且整个东西是彩色的。所有这些将在适当的时候加以解释，但就目前来说，它确是说明光并不总是沿着直线行进的一个演示实验，而且是一个很容易做的演示实验。

\section{最短时间原理是怎样起作用的}

最后，从我们现在相信是正确的、量子动力学上精确的观点出发，对实际发生的是什么以及整个事情是怎样起作用的提供一个粗略的概念，当然只能作定性的描述。在图\ref{figure:最短时间原理的说明}中，当我们随着光从$A$到$B$行进时，我们发现光似乎根本不具有波的形式。相反，光线倒是有点像由光子组成的，如果我们用一个光子计数器，它实际上会在其中产生咔嗒声。光的亮度与每秒钟进入计数器的平均光子数成正比，而我们所计算的则是光子从$A$（比如说碰到镜面后）到达$B$的机会。这种机会所遵循的是下述很奇怪的规律。取任一路径，并找出其相应的时间；然后写一复数或画一小的复矢量$\rho e^{i\theta}$，令其角度$\theta$正比于时间。复矢量每秒的旋转周数就是光的频率。再取另一路径，比如说它具有不同的时间，则其对应的矢量就转过不同的角度——角度总是与时间成正比的。取所有可取的路径，并为每一条路径加上一个小矢量；那么答案就是，光子到达的机会与从始端到末端的总矢量的长度平方成正比！


现在我们来说明对镜面来说这个结果如何隐含了最短时间原理。考虑图\ref{figure:最短时间原理的说明}中所有
\begin{wrapfigure}{r}{0.35\textwidth}
    \centering
    \includegraphics[width=0.34\textwidth]{Chapter26/许多邻近路径的概率振幅总和}
    \caption{许多邻近路径的概率振幅总和}
    \label{figure:许多邻近路径的概率振幅总和}
\end{wrapfigure}
的光线和所有可能的路径$ADB$，$AEB$，$ACB$等等。路径$ADB$提供了某一小的贡献，但下一个路径$AEB$所花的时间就完全不同，故其角度$\theta$也完全不同。设$C$点对应于时间的极小值，在它附近改变路径时，时间并不改变。因而，起初一会儿时间在改变，但当接近$C$点时，时间的变化就开始变得越来越小（图\ref{figure:许多邻近路径的概率振幅总和}）。所以在靠近$C$点的片刻所添加的箭头几乎都达到完全相同的角度。然后时间又开始逐渐增加，相角又朝相反方向旋转，等等。最后，得到一个收紧的结。总的概率就是从一端到另一端距离的平方。几乎所有对于概率的积累都发生在所有箭头沿同一方向（即同相位）的区域。所有来自那些当改变路径时具有十分不同时间的路径的贡献，由于指向不同方向而相互抵消了。这就是为什么当我们遮去镜子的边缘部分时，它的反射几乎与以前相同的缘故，因为我们在这里所作的一切仅仅是抹去了位于图上的螺旋末端的一部分而已，它只引起光很少的改变。这就是最终的（其到达的概率取决于复矢量的累加）光子图像与最短时间原理之间的关系。